% Partie cours issue de 7-IntroductionSA2022.docx.
Objectifs du cours

A la fin du cours, vous devez être capables de~:

\begin{itemize}
\item
  \textbf{connaître} la différence entre un système bouclé et un système
  en boucle ouverte~;
\item
  \textbf{connaître la structure des systèmes asservis (chaîne directe,
  chaîne de retour,\ldots)}
\item
  \textbf{connaître} les deux types de systèmes asservis~;
\item
  \textbf{mettre} un système asservi sous forme de schéma-blocs~;
\item
  \textbf{connaître et caractériser} les performances des systèmes
  asservis.
\end{itemize}

Savoirs

\textbf{Je connais~:}

\begin{longtable}[]{@{}
  >{\raggedright\arraybackslash}p{(\columnwidth - 2\tabcolsep) * \real{0.9045}}
  >{\raggedright\arraybackslash}p{(\columnwidth - 2\tabcolsep) * \real{0.0955}}@{}}
\toprule\noalign{}
\begin{minipage}[b]{\linewidth}\raggedright
\begin{itemize}
\item
  les différents types de systèmes automatisés~;
\end{itemize}
\end{minipage} & \begin{minipage}[b]{\linewidth}\raggedright
$\bigcirc$
\end{minipage} \\
\midrule\noalign{}
\endhead
\bottomrule\noalign{}
\endlastfoot
\begin{minipage}[t]{\linewidth}\raggedright
\begin{itemize}
\item
  les différents types de systèmes asservis~;
\end{itemize}
\end{minipage} & $\bigcirc$ \\
\begin{minipage}[t]{\linewidth}\raggedright
\begin{itemize}
\item
  la structure d'un système asservi~;
\end{itemize}
\end{minipage} & $\bigcirc$ \\
\begin{minipage}[t]{\linewidth}\raggedright
\begin{itemize}
\item
  les critères de performances~;
\end{itemize}
\end{minipage} & $\bigcirc$ \\
\begin{minipage}[t]{\linewidth}\raggedright
\begin{itemize}
\item
  les entrées tests.
\end{itemize}
\end{minipage} & $\bigcirc$ \\
\end{longtable}

Savoir-faire

\textbf{Je sais~:}

\begin{longtable}[]{@{}
  >{\raggedright\arraybackslash}p{(\columnwidth - 2\tabcolsep) * \real{0.9045}}
  >{\raggedright\arraybackslash}p{(\columnwidth - 2\tabcolsep) * \real{0.0955}}@{}}
\toprule\noalign{}
\begin{minipage}[b]{\linewidth}\raggedright
\begin{itemize}
\item
  reconnaître un système asservi~;
\end{itemize}
\end{minipage} & \begin{minipage}[b]{\linewidth}\raggedright
$\bigcirc$
\end{minipage} \\
\midrule\noalign{}
\endhead
\bottomrule\noalign{}
\endlastfoot
\begin{minipage}[t]{\linewidth}\raggedright
\begin{itemize}
\item
  caractériser les performances d'un système asservi.
\end{itemize}
\end{minipage} & $\bigcirc$ \\
\end{longtable}

\section{MISE EN SITUATION}\label{mise-en-situation}

\subsection[Cas n$^\circ$1~: Etuve
thermique]{Cas n$^\circ$1~: Etuve thermique}\label{cas-n1-etuve-thermique}

L'étuve thermique est utilisée pour la stérilisation des ustensiles de
laboratoire (tubes à essai, boîtes de Pétri, tubes à culture et
bouchons, pipettes Pasteur et récipients divers). Elle permet un réglage
de température de la température ambiante jusqu\textquotesingle à 200
$^\circ$C.

Pour supprimer les micro-organismes, il est nécessaire de chauffer
pendant 2 heures à une température constante de 150$^\circ$C .

\SAimage{images/image6.png}{78mm}{52mm}Une
résistance électrique alimentée par une tension \(U(t)\) assure le
chauffage à l\textquotesingle intérieur de l\textquotesingle étuve. Une
sonde de température Pt100, associée à un conditionneur permet de
mesurer la température \(\theta(t)\) dans l\textquotesingle étuve.

Un régulateur, composé d'un comparateur et d'un correcteur, élabore la
tension de commande \(U_{com}(t)\) du gradateur à partir de la
différence entre la température de consigne \(\theta_{cons}(t)\) et la
température mesurée \(\theta_{mes}(t)\).

\textbf{Comment garantir que la température restera toujours constante à
150$^\circ$C pendant 2h~?}

\subsection[Cas n$^\circ$2~: Robot Da
Vinci]{Cas n$^\circ$2~: Robot Da Vinci}\label{cas-n2-robot-da-vinci}

Un Da Vinci est un robot médical, et plus précisément une machine
dirigée par un chirurgien pour réaliser des opérations, principalement
au niveau de l\textquotesingle abdomen.

Devenue un standard aux États-Unis, la chirurgie robotique se diffuse en
France depuis quelques années. Outre-Atlantique, plus de 80 \% des cas
de cancer de la prostate sont opérés à l\textquotesingle aide du robot
Da Vinci. Ses avantages ? Opérations moins invasives, récupération du
patient plus rapide, cicatrices plus petites, voire invisibles,
complications et infections postopératoires diminuées, pertes sanguines
réduites durant l\textquotesingle opération, risque diminué de séquelles
liées à l\textquotesingle opération. Tous ces facteurs améliorent la
prise en charge du patient et diminuent la durée de son hospitalisation.

Le Da Vinci est composé de deux parties. La première se situe au-dessus
de la personne à opérer et comporte trois (dans sa première version) ou
quatre (à partir du modèle Da Vinci S) bras manipulateurs. Un bras tient
une caméra endoscopique, les autres tiennent des instruments
chirurgicaux tels qu\textquotesingle un bistouri, ou plus précisément un
électrobistouri qui découpe les tissus à l\textquotesingle aide
d\textquotesingle un courant électrique.

La seconde est située à quelques mètres de la première, et comporte un
siège sur lequel s\textquotesingle assied le chirurgien, deux écrans
devant lesquels ce dernier vient placer ses yeux et qui retransmettent
en direct la vue en 3D de la caméra endoscopique située sur la première
partie, et deux manettes pour contrôler les instruments chirurgicaux
situés sur la première partie afin de retranscrire fidèlement les
mouvements complexes des mains du chirurgien à l\textquotesingle aide de
pinces articulées et miniaturisées.

\textbf{Comment être sûr que le robot effectue les mouvements réellement
demandés par le chirurgien~?}

\subsection{Cas n$^\circ$3~: Régulateur de vitesse
Renault}\label{cas-n3-ruxe9gulateur-de-vitesse-renault}

\SAimage{images/image8.jpeg}{64mm}{45mm}Le
conducteur d\textquotesingle une Laguna a parcouru 47 km sur
l\textquotesingle autoroute A 13 en Normandie au volant de sa voiture
folle, dont le régulateur de vitesse est resté bloqué. Ce jour-là, au
volant de sa Renault Laguna, le jeune homme roule sur une portion
limitée à 90 km/h qui mène à l\textquotesingle autoroute A 13. «
J\textquotesingle ai bloqué le régulateur à 87 km/h afin de circuler en
toute sécurité. Mais quand il a fallu que je freine parce
qu\textquotesingle un véhicule ralentissait, les freins ne répondaient
plus du tout. J\textquotesingle ai donc été obligé de dépasser
brusquement cette voiture. Je ne pouvais plus rien faire, ni accélérer,
ni ralentir, ni débrayer, ni passer une vitesse, ni enlever la carte
électronique. La voiture était incontrôlable, le régulateur de vitesse
était bloqué. J\textquotesingle ai paniqué. »

Pour maitriser parfaitement une grandeur physique (position, vitesse,
température, \ldots) et ce quelles que soient les perturbations, il faut
asservir cette grandeur~!


\section{INTRODUCTION ET DÉFINITIONS}\label{introduction-et-definitions}

\subsection{Système automatisé}\label{systuxe8me-automatisuxe9}

Les échanges entre la chaîne d'information et la chaîne de puissance
permettent à certaines activités d'un système d'être réalisées de
manière relativement autonome. On parle de \textbf{système automatisé}.
Un \textbf{système automatisé} est un système réalisant des
\textbf{opérations sans intervention de l'homme}. Les systèmes
automatisés sont utilisés pour réaliser des tâches :

\begin{itemize}
\item
  \textbf{complexes} ou \textbf{dangereuses} pour
  l\textquotesingle homme~;
\end{itemize}

\SAimagepair{images/image11.png}{0.24\linewidth}{images/image12.png}{0.46\linewidth}

Correction de la vue par laser

Robot serpent Tepco d'inspection du cœur d'une centrale nucléaire

\begin{itemize}
\item
  \textbf{répétitives} ou \textbf{pénibles}.
\end{itemize}

\SAimagepair{images/image13.jpeg}{0.44\linewidth}{images/image14.jpeg}{0.44\linewidth}

Voiture autonome Tesla

Robots guides d'un complexe financier Santander de Madrid

\SAimage{images/image150.pdf}{100mm}{55mm}

L'intervention de l'homme est alors limitée à la \textbf{programmation}
de l'unité de commande, la mise en marche et les \textbf{réglages} de
certains \textbf{paramètres}.

À partir du \textbf{modèle du système} étudié, sous forme d'équations,
il s'agira \textbf{d'évaluer ses performances} dans le but d'anticiper
son comportement réel et vérifier ainsi s'il est capable de répondre aux
\textbf{exigences} de son \textbf{cahier des charges}.

Notons enfin que, bien que de nombreux systèmes utilisent actuellement
des échanges et un traitement des informations sous forme numérique, on
limitera nos études uniquement aux cas des systèmes continus. Dans un
système \textbf{continu}, les \textbf{grandeurs} \textbf{d'entrée} et de
\textbf{sortie} évoluent de manière \textbf{continue} en fonction du
\textbf{temps}, il s'agit donc de grandeurs analogiques.

\subsection{Notion de système bouclé / boucle
ouverte}\label{notion-de-systuxe8me-boucluxe9-boucle-ouverte}

\textbf{Système}

\textbf{Industriel}

\textbf{Ordre ou Sollicitation ou Cause ou Consigne}

\textbf{Réponse ou Effet}

\emph{\textbf{Perturbations}}

Un \textbf{système non bouclé (ou en boucle ouverte)} est un système qui
\textbf{ne contrôle pas la manière dont l'ordre a été exécuté}. Il
\textbf{ne prend pas en compte la réaction du système}.

Un événement extérieur (\textbf{perturbation}) peut alors modifier la
\textbf{réponse} (ou grandeur de sortie) attendue pour le système.

\textbf{Etuve thermique}

\textbf{Consigne de température}

\textbf{Chauffage de l'étuve}

\emph{\textbf{Variations de la température ambiante, ouverture de la
porte}}

\subparagraph*{Exemple de système en boucle ouverte }\label{exemple-de-systuxe8me-en-boucle-ouverte}
\addcontentsline{toc}{subparagraph}{Exemple de système en boucle ouverte
}

Un conducteur de voiture qui conduit sur une ligne droite ferme les yeux
et maintient le volant. La voiture va continuer à avancer à une vitesse
constante. Si un piéton traverse ou qu'une voiture s'arrête
(perturbations) le conducteur n'aura pas le temps de les éviter et ce
sera l'accident. Pour les éviter le conducteur a besoin d'observer en
permanence la route.

\textbf{Autres exemples :} chauffage d'immeuble avec fenêtre ouverte,
arrosage automatique sans contrôle d'humidité, etc...

De nombreux systèmes nécessitent que \textbf{la grandeur de sortie
(réponse) ne varie pas} (température, la vitesse d'un moteur, la
position d'un bras de robot, le cap d'un bateau, ... ) \textbf{quels que
soient les phénomènes extérieurs ou perturbations} (ouverture de la
fenêtre dans un système de chauffage, variation du couple résistant ...)
qui peuvent influer sur ces grandeurs.

Pour pallier à ce problème, il faut \textbf{observer l'état du système
pour le modifier et obtenir la réponse souhaitée quelles que soient les
perturbations.}

\textbf{Capteur}

Réponse

Consigne

Mesure

+

-

Ecart

\textbf{Système}

\textbf{industriel}

\emph{\textbf{Perturbations}}

Il est alors réalisé une \textbf{rétroaction} \textbf{ou contre-réaction
(feedback} en anglais), pour \textbf{observer l'état du système afin de
le modifier et d'obtenir la réponse souhaitée quelles que soient les
perturbations.}

Un \textbf{système bouclé (ou en boucle fermée)} est un système dont la
\textbf{grandeur de sortie est mesurée puis comparée à la valeur de
consigne} (rétroaction de la sortie sur l'entrée). Cette différence est
appelée \textbf{l'écart}.

Le but des systèmes asservis est \textbf{d'annuler en permanence l'écart
entre la grandeur de sortie et l'entrée (consigne)}. La sortie est alors
\textbf{asservie} à l'entrée.

Un système asservi \textbf{comporte toujours une boucle de retour}
(contre-réaction) comprenant un capteur.

\subsection{Structure générale d'un système
asservi}\label{structure-guxe9nuxe9rale-dun-systuxe8me-asservi}

La structure complète d'un système asservi peut se représenter par le
schéma-bloc fonctionnel suivant.

\SASchemaBoucleGenerale

Il est constitué de deux chaînes :

\begin{itemize}
\item
  La \textbf{chaîne directe} (ou d'action) assurant les fonctions de
  commande et puissance. On y retrouve les constituants de la chaine
  d'énergie.
\item
  La \textbf{chaîne de retour} (ou de réaction) assurant la fonction de
  mesure. On y retrouve les constituants de la chaine d'information.
\end{itemize}

Le \textbf{régulateur} élabore le signal de commande à partir de
\textbf{l'écart $\varepsilon$} constaté entre la \textbf{consigne} et la
\textbf{mesure} issue du \textbf{capteur}. Il comporte~:

\begin{itemize}
\item
  un \textbf{comparateur~} élaborant le signal d'écart \textbf{$\varepsilon$}~;
\item
  un \textbf{correcteur~}modifiant l'allure de la commande pour
  améliorer les performances du système.
\end{itemize}

L'ensemble «pré-actionneur, actionneur, adaptateur» réalise la commande
de l'effecteur à partir du signal de commande élaboré par le régulateur.

\subsection{Types de systèmes
asservis}\label{types-de-systuxe8mes-asservis}

Il existe deux types de systèmes asservis :

\begin{itemize}
\item
  Les \textbf{systèmes régulateurs~}qui maintiennent \textbf{la grandeur
  de sortie constante quelles que soient les perturbations} ;

  \begin{itemize}
  \item
    \ul{Ex:} Régulation de vitesse, de température, \ldots{}
  \end{itemize}
\item
  Les \textbf{systèmes suiveurs} dont \textbf{la grandeur de sortie suit
  la loi de variation donnée en consigne} (entrée).

  \begin{itemize}
  \item
    \ul{Ex:} Radar de poursuite, anti-missiles, machine d'usinage, robot
    suiveur de trajectoire \ldots{}
  \end{itemize}
\end{itemize}

\SAimage{images/image16.jpeg}{25mm}{32mm}\textbf{Régulateur
du niveau d'eau d'une piscine :} Système de régulation du niveau d'eau
pour compenser les éclaboussures et l'évaporation.

L'étude des systèmes asservis permet de répondre aux questions
suivantes:

\begin{itemize}
\item
  Comment va réagir le système à une consigne donnée (\textbf{Aspect
  prédictif})
\item
  Quel type de système fournit telle réponse à telle sollicitation?
  (\textbf{Aspect caractérisation, identification})
\item
  Comment \textbf{adapter ou régler un système} pour que ses
  caractéristiques correspondent à celles définies dans le cahier des
  charges~?
\end{itemize}

\subsection{Représentation par
schéma-blocs}\label{repruxe9sentation-par-schuxe9ma-blocs}

En automatique, les systèmes sont représentés par un ensemble de
\textbf{schéma-blocs fonctionnels} : c'est la représentation
fonctionnelle.

Les 3 éléments de base sont :

\begin{itemize}
\item le \textbf{bloc}, qui représente un constituant ou une fonction de transfert ;
\item le \textbf{comparateur (ou sommateur)}, qui additionne ou soustrait plusieurs signaux ;
\item le \textbf{point de prélèvement}, qui duplique un signal sans le modifier.
\end{itemize}

\SASchemaElementsBase

\textbf{Exemple :} un moteur peut être représenté par un bloc ayant une tension de commande en entrée et une vitesse de rotation en sortie.

\SASchemaMoteur

La représentation par schéma-blocs d'un système peut être obtenue à
partir de différentes formes:

\begin{itemize}
\item
  à partir des chaînes fonctionnelles du système ;
\end{itemize}

\SAChaineFonctionnelle

\begin{itemize}
\item
  à partir d'un texte ou d'un schéma décrivant la structure du système ;
\end{itemize}

\SAChaineAsservie

\begin{longtable}[]{@{}
  >{\raggedright\arraybackslash}p{(\columnwidth - 2\tabcolsep) * \real{0.1227}}
  >{\raggedright\arraybackslash}p{(\columnwidth - 2\tabcolsep) * \real{0.8773}}@{}}
\toprule\noalign{}
\begin{minipage}[b]{\linewidth}\raggedright
\begin{quote}
\SAimage{images/image18.png}{13mm}{14mm}
\end{quote}
\end{minipage} & \begin{minipage}[b]{\linewidth}\raggedright
\SAimage{images/image6.png}{78mm}{52mm}\SAimage{images/image5.png}{30mm}{30mm}\textbf{Etuve
thermique}

Une résistance électrique alimentée par une tension \(U(t)\) assure le
chauffage à l\textquotesingle intérieur de l\textquotesingle étuve. Une
sonde de température Pt100, associée à un conditionneur permet de
mesurer la température \(\theta(t)\) dans l\textquotesingle étuve.

Un régulateur, composé d'un comparateur et d'un correcteur, élabore la
tension de commande \(U_{com}(t)\) du gradateur à partir de la
différence entre la température de consigne \(\theta_{cons}(t)\) et la
température mesurée \(\theta_{mes}(t)\).

\textbf{Compléter la représentation, sous forme de schéma-bloc
fonctionnel, de l'étuve.}

\SASchemaEtuve
\end{minipage} \\
\midrule\noalign{}
\endhead
\bottomrule\noalign{}
\endlastfoot
\end{longtable}

\section{Performances d'un système
asservi}\label{performances-dun-systuxe8me-asservi}

Un système est asservi pour que son fonctionnement soit conforme à des
attentes définies dans un cahier des charges. Ces attentes sont
déclinées en termes de \textbf{performances}.

L'objectif est d'obtenir, à l'aide de réglages, le \textbf{meilleur
compromis} entre les différents critères de performances globales du
système et de valider les critères du CdCF.

\subsection{Critères de performance d'un système
asservi}\label{crituxe8res-de-performance-dun-systuxe8me-asservi}

Les performances sont des caractéristiques propres à un système ou liées
à son comportement lorsqu'il est sollicité de façon particulière. Nous
allons voir les critères de performances suivants :

\begin{itemize}
\item
  la \textbf{stabilité}, le système est-il stable? \textbf{C'est le
  critère le plus important} !! C'est la propriété de convergence
  temporelle vers un état d'équilibre ;
\item
  la \textbf{rapidité}, caractérisant le dynamisme des systèmes
  stables~;
\item
  la \textbf{précision}, aptitude des systèmes asservis stables à
  atteindre la valeur attendue.
\item
  Le \textbf{dépassement} : Est-ce que la réponse dépasse sa valeur
  finale transitoirement~?
\end{itemize}

Pour \textbf{évaluer} ces différents critères, il est nécessaire de
\textbf{procéder à des tests soit sur le système réel, soit en
simulation} en choisissant les entrées qui peuvent être appliquées.

La précision, la rapidité et le dépassement \textbf{peuvent être évalués
à partir des réponses aux signaux de test} (réponse temporelle). La
stabilité est évaluée à partir de la \textbf{réponse harmonique (entrée
sinusoïdale).}

Afin d'évaluer le comportement (puis de le corriger) d'un système
asservi, il y a plusieurs possibilités :

\begin{itemize}
\item
  Réaliser directement des mesures sur le système réel en appliquant des
  signaux tests~;
\item
  Réaliser une simulation / calculs à partir du modèle du système en
  appliquant des signaux tests.
\end{itemize}

\textbf{Simulation avec le modèle du système}

\textbf{Réponse simulée}

\textbf{Réponse expérimentale}

\emph{\textbf{Comparaison}}

\emph{\textbf{Vérification des performances}}

\emph{\textbf{Vérification des performances}}

\textbf{Choix de l'entrée (consigne)}

\textbf{Mesures sur le système réel}

\textbf{Cahier des charges}

\subsection{Signaux tests (ou
canoniques)}\label{signaux-tests-ou-canoniques}

Les performances d'un système sont évaluées expérimentalement, ou
prédites par simulation, lorsque le système est soumis à des
\textbf{signaux tests e(t) en entrée} appelés \textbf{signaux
canoniques}.

Système

e(t)

consigne

s(t)

réponse

Ces signaux tests envoyés en entrée sont~:

\begin{longtable}[]{@{}
  >{\raggedright\arraybackslash}p{(\columnwidth - 2\tabcolsep) * \real{0.5000}}
  >{\raggedright\arraybackslash}p{(\columnwidth - 2\tabcolsep) * \real{0.5000}}@{}}
\toprule\noalign{}
\begin{minipage}[b]{\linewidth}\raggedright
\textbf{L'impulsion (dite de Dirac)}

\textbf{Réponse impulsionnelle}

\[\varepsilon\]

\[\frac{1}{\varepsilon}\]

\emph{t}

\emph{$\delta$(t)}

\emph{avec}\(\ \varepsilon \rightarrow 0\)

Exemple~: Airbag
\end{minipage} & \begin{minipage}[b]{\linewidth}\raggedright
\textbf{L'échelon}

\textbf{Réponse indicielle}

\SAimage{images/image22.png}{65mm}{65mm}

\emph{u(t)}

\[1\]

\[0\]

\emph{t}

Exemple~: Thermostat
\end{minipage} \\
\midrule\noalign{}
\endhead
\bottomrule\noalign{}
\endlastfoot
\textbf{La rampe}

\SAimage{images/image23.png}{38mm}{46mm}

\emph{t u(t)}

\emph{t}

Exemple~: Tracker solaire & \textbf{La sinusoïde}

\textbf{Réponse harmonique}

\SAimage{images/image24.png}{55mm}{34mm}

\emph{sin(t) u(t)}

\emph{t}

Exemple~: Suspension \\
\end{longtable}

Par convention, tous ces signaux sont nuls pour t négatif~:
\(\mathbf{e(t) = 0\ }\)\textbf{pour} \(\mathbf{t < 0}\)

Notons toutefois que les performances sont \textbf{intrinsèques} au
système et \textbf{ne dépendent pas du type d'entrée appliqué.}

La rampe est la \textbf{primitive} de l\textquotesingle échelon qui est
la \textbf{primitive} de l\textquotesingle impulsion. On pourrait le
généraliser par intégration successive.

L\textquotesingle entrée sinusoïdale se généralise en entrée harmonique.

\textbf{Les systèmes asservis étudient uniquement les variations autour
d'un point de fonctionnement (ou point de repos).}

\subparagraph*{Exemple de choix des
entrées}\label{exemple-de-choix-des-entruxe9es}
\addcontentsline{toc}{subparagraph}{Exemple de choix des entrées}

\begin{longtable}[]{@{}
  >{\raggedright\arraybackslash}p{(\columnwidth - 2\tabcolsep) * \real{0.1227}}
  >{\raggedright\arraybackslash}p{(\columnwidth - 2\tabcolsep) * \real{0.8773}}@{}}
\toprule\noalign{}
\begin{minipage}[b]{\linewidth}\raggedright
\begin{quote}
\SAimage{images/image18.png}{13mm}{14mm}
\end{quote}
\end{minipage} & \begin{minipage}[b]{\linewidth}\raggedright
\SAimage{images/image5.png}{30mm}{30mm}\textbf{Etuve
thermique}

On souhaite évaluer les performances (rapidité, précision, dépassement)
de l'étuve en prenant le cas suivant :

Afin de stériliser les ustensiles de laboratoire, l'utilisateur
programme la température de l'étuve à 150$^\circ$C. La température ambiante (et
à l'intérieur de l'étuve \(\theta(t)\)) est de 25$^\circ$C.

\textbf{Etuve thermique}

\textbf{$\theta$(t)}

\textbf{$\theta$\textsubscript{cons}(t)}

\textbf{Quelle entrée}
\(\mathbf{\theta}_{\mathbf{cons}}\mathbf{(}\mathbf{t}\mathbf{)}\)
\textbf{doit-on utiliser pour évaluer ces critères de performances ?}

\[150\]

\[25\]

\[\theta(t)\]

\[t\]

\[t_{init}\]

\[\Delta\theta\]

La consigne est constante (150$^\circ$C), l'entrée sera donc de type échelon.
La variation de température autour du point de repos
(\(\theta(t)\)=25$^\circ$C) est \(\Delta\theta\) =125$^\circ$C.

Les systèmes asservis étudiant uniquement les variations autour d'un
point de fonctionnement, il est possible d'évaluer ces critères en
simulation avec un échelon d'amplitude 125$^\circ$. \textbf{La forme de la
réponse sera la même}.

\[125\]

\[0\]

\[\theta(t)\]

\[t\]

\[t_{init}\]

\[\Delta\theta\]
\end{minipage} \\
\midrule\noalign{}
\endhead
\bottomrule\noalign{}
\endlastfoot
\end{longtable}

\subsection{Consigne et réponse}\label{consigne-et-ruxe9ponse}

Les valeurs finales des grandeurs d'entrée (consigne) et de sortie
(réponse), correspondent aux valeurs de l'entrée et de la sortie du
système en régime permanent, soit pour un temps suffisamment grand.
Elles sont notées \(e_{\infty}\) et \(s_{\infty}\) avec~:

\(e_{\infty} = \lim_{t \rightarrow + \infty}{e(t)}\) et
\(s_{\infty} = \lim_{t \rightarrow + \infty}{s(t)}\)

La variation totale de la grandeur d'entrée est
\(\Delta E = e_{\infty} - e\left( t_{init} \right)\)

La variation totale de la grandeur de sortie est
\(\Delta S = s_{\infty} - s(t_{init})\)

\[s(t)\]

\[s_{\infty}
\]

\[s(t_{init})\]

\[t\]

\[t_{init}\]

\[\Delta S\]

\[e_{\infty}\]

\[
{e\left( t_{init} \right)}\]

\[e(t)\]

\[t\]

\[t_{init}\]

\[\Delta E\]

\subsection{Stabilité des
systèmes}\label{stabilituxe9-des-systuxe8mes}

La \textbf{stabilité est la performance qu'il faut évaluer en premier}
car un système instable est inutilisable.

\textbf{La sortie diverge}

\textbf{Système Instable}

t

0

s(t)

S\textsubscript{$\infty$}

\textbf{Système Stable}

\textbf{La sortie converge}

t

0

s(t)

S\textsubscript{$\infty$}

La stabilité est évaluée à partir de la \textbf{réponse harmonique} ou
fréquentielle et peut s'observer à partir de la réponse indicielle.

Un système est \textbf{stable} si, \textbf{pour une entrée en échelon},
la grandeur de sortie \textbf{converge} vers une \textbf{valeur finale
constante}.

(entrée bornée -- sortie bornée EB-SB).

Un système réel instable oscille jusqu'à sa destruction. Ces
oscillations sont dans le cas général \textbf{limitées par les
différentes saturations} (électroniques ou mécaniques\ldots). Ces
limitations physiques peuvent laisser croire que le système est stable.

\subsection{Rapidité des systèmes}\label{rapidituxe9-des-systuxe8mes}

Pour certains systèmes, il est impératif que le temps de réponse du
système soit faible. À condition que le \textbf{système soit stable}, la
\textbf{rapidité} est caractérisée, pour une \textbf{entrée en échelon
(réponse indicielle)}, par la durée que met le système pour que la
sortie soit suffisamment proche de sa valeur finale et ne plus s'en
éloigner.

Par convention, le critère caractérisant la \textbf{rapidité} d'un
système stable est le \textbf{temps de réponse à 5\%} noté
\(t_{r5\%}\)\emph{.} On définit le \textbf{tube des 5\%} par
l\textquotesingle intervalle~:

\(\left\lbrack s_{\infty} - 0,05\ \mathrm{\Delta}S\ ,\ s_{\infty} + 0,05\ \mathrm{\Delta}S \right\rbrack\)

\subparagraph*{Réponse d'un système pour une entrée test en
échelon}\label{ruxe9ponse-dun-systuxe8me-pour-une-entruxe9e-test-en-uxe9chelon}
\addcontentsline{toc}{subparagraph}{Réponse d'un système pour une entrée
test en échelon}

\[s_{\infty}\]

\[t_{init}\]

\[t\]

\[s(t)\]

\[t_{r5\%}\]

\[s_{\infty} - 0,05\ \mathrm{\Delta}s_{\infty}\ \]

\[\Delta S\]

\[tube\ des\ 5\%\]

\[bande\ des\ 5\%\]

\[s_{\infty} + 0,05\ \mathrm{\Delta}s_{\infty}\]

Le \textbf{temps de réponse à 5\%} est la \textbf{durée} mise par la
grandeur de \textbf{sortie} pour rentrer dans le tube des 5\% et ne
\textbf{plus en sortir}.

\paragraph*{On pourrait utiliser d'autres critères, comme le temps de montée (premier instant auquel la réponse atteint la valeur finale) ou le temps de réponse à 1\%. Cette valeur de 5\% vient de la réponse temporelle d'un premier ordre. On ne dira jamais que le temps de réponse est rapide, bon... il faut comparer. }\label{on-pourrait-utiliser-dautres-crituxe8res-comme-le-temps-de-montuxe9e-premier-instant-auquel-la-ruxe9ponse-atteint-la-valeur-finale-ou-le-temps-de-ruxe9ponse-uxe0-1.-cette-valeur-de-5-vient-de-la-ruxe9ponse-temporelle-dun-premier-ordre.-on-ne-dira-jamais-que-le-temps-de-ruxe9ponse-est-rapide-bon...-il-faut-comparer.}
\addcontentsline{toc}{paragraph}{On pourrait utiliser d'autres critères,
comme le temps de montée (premier instant auquel la réponse atteint la
valeur finale) ou le temps de réponse à 1\%. Cette valeur de 5\% vient
de la réponse temporelle d'un premier ordre. On ne dira jamais que le
temps de réponse est rapide, bon... il faut comparer. }

\subsubsection{ Il faut construire le tube des 5\% autour de la valeur finale de la sortie et non pas autour de la valeur de la grandeur d'entrée~!}\label{il-faut-construire-le-tube-des-5-autour-de-la-valeur-finale-de-la-sortie-et-non-pas-autour-de-la-valeur-de-la-grandeur-dentruxe9e}

\subsubsection*{}\label{section}
\addcontentsline{toc}{subsubsection}{}

\begin{longtable}[]{@{}
  >{\raggedright\arraybackslash}p{(\columnwidth - 2\tabcolsep) * \real{0.1227}}
  >{\raggedright\arraybackslash}p{(\columnwidth - 2\tabcolsep) * \real{0.8773}}@{}}
\toprule\noalign{}
\begin{minipage}[b]{\linewidth}\raggedright
\begin{quote}
\SAimage{images/image18.png}{13mm}{14mm}
\end{quote}
\end{minipage} & \begin{minipage}[b]{\linewidth}\raggedright
\SAimage{images/image5.png}{30mm}{30mm}\textbf{Etuve
thermique}

Le robot Da Vinci est un robot médical dirigé par un chirurgien pour
réaliser des opérations, principalement au niveau de l'abdomen.
L'opération se déroule à distance à l'aide d'un écran et de dispositifs
haptiques. Les mouvements sont tous asservis et requièrent un réglage
fin.

On souhaite évaluer les performances (rapidité, précision, dépassement)
de l'étuve en prenant le cas suivant :

\textbf{Etuve thermique}

\textbf{$\theta$(t)}

\textbf{$\theta$\textsubscript{cons}(t)}

Afin de stériliser les ustensiles de laboratoire, l'utilisateur
programme la température de l'étuve à 150$^\circ$C. La température ambiante (et
à l'intérieur de l'étuve \(\theta(t)\)) est de 25$^\circ$C.

\textbf{Donner le temps de réponse du système non corrigé.}

\SAimage{images/image27.png}{88mm}{72mm}

\textbf{Donner le temps de réponse du système corrigé. Conclure.}

\SAimage{images/image28.png}{88mm}{72mm}
\end{minipage} \\
\midrule\noalign{}
\endhead
\bottomrule\noalign{}
\endlastfoot
\end{longtable}

\begin{longtable}[]{@{}
  >{\raggedright\arraybackslash}p{(\columnwidth - 2\tabcolsep) * \real{0.1227}}
  >{\raggedright\arraybackslash}p{(\columnwidth - 2\tabcolsep) * \real{0.8773}}@{}}
\toprule\noalign{}
\begin{minipage}[b]{\linewidth}\raggedright
\begin{quote}
\SAimage{images/image18.png}{13mm}{14mm}
\end{quote}
\end{minipage} & \begin{minipage}[b]{\linewidth}\raggedright
\SAimage{images/image29.png}{68mm}{38mm}\textbf{Robot
chirurgical Da Vinci}

\textbf{Déterminer le temps de réponse à 5\%.}

x(t) {[}mm{]}

\SAimage{images/image30.pdf}{90mm}{75mm}

t {[}ms{]}

\begin{quote}
On lit graphiquement

\[x_{\infty} = 8mm\]

\[\Delta X = 8 - 5 = 3\ mm\]

Le "tube des 5\%" est~:

{[}8-0,05x3~; 8+0,05x3{]} = {[}7,85~;~8,15{]}

On lit graphiquement :

\[t_{r5\%} = 100 - 20 = 80\ ms\]
\end{quote}
\end{minipage} \\
\midrule\noalign{}
\endhead
\bottomrule\noalign{}
\endlastfoot
\end{longtable}

\begin{longtable}[]{@{}
  >{\raggedright\arraybackslash}p{(\columnwidth - 2\tabcolsep) * \real{0.1227}}
  >{\raggedright\arraybackslash}p{(\columnwidth - 2\tabcolsep) * \real{0.8773}}@{}}
\toprule\noalign{}
\begin{minipage}[b]{\linewidth}\raggedright
\begin{quote}
\SAimage{images/image18.png}{13mm}{14mm}
\end{quote}
\end{minipage} & \begin{minipage}[b]{\linewidth}\raggedright
\SAimage{images/image31.jpeg}{28mm}{58mm}\textbf{Stabilisateur}

Un essai indiciel sur le stabilisateur pour un déplacement de 2$^\circ$ est
donné ci-dessous.

\SAimage{images/image32.pdf}{90mm}{75mm}

\textbf{Déterminer le temps de réponse à 5\%.}
\end{minipage} \\
\midrule\noalign{}
\endhead
\bottomrule\noalign{}
\endlastfoot
\end{longtable}

\begin{longtable}[]{@{}
  >{\raggedright\arraybackslash}p{(\columnwidth - 2\tabcolsep) * \real{0.1227}}
  >{\raggedright\arraybackslash}p{(\columnwidth - 2\tabcolsep) * \real{0.8773}}@{}}
\toprule\noalign{}
\begin{minipage}[b]{\linewidth}\raggedright
\begin{quote}
\SAimage{images/image18.png}{13mm}{14mm}
\end{quote}
\end{minipage} & \begin{minipage}[b]{\linewidth}\raggedright
\textbf{Stabilisateur essai 2}

Un essai indiciel sur le stabilisateur pour un déplacement de 2$^\circ$ est
donné ci-dessous.

\SAimage{images/image33.png}{90mm}{74mm}

\textbf{Déterminer le temps de réponse à 5\%.}
\end{minipage} \\
\midrule\noalign{}
\endhead
\bottomrule\noalign{}
\endlastfoot
\end{longtable}

\subsection{Précision des systèmes
asservis}\label{pruxe9cision-des-systuxe8mes-asservis}

Pour certains systèmes, il est impératif que la précision du système
soit excellente. La \textbf{précision s'évalue à partir de la
différence} entre l'entrée (consigne ou valeur souhaitée) et la sortie
(réponse effective) \textbf{en régime permanent} (quand l'écart entre la
consigne et la réponse n'évolue plus).

Elle est \textbf{évaluée pour une entrée donnée} (échelon, rampe,...).

s(t)

t

0

\textbf{Réponse à un échelon (indicielle)}

\textbf{$\varepsilon$\textsubscript{s}}

\textbf{Ecart statique ou écart de position $\varepsilon$\textsubscript{s}}

S\textsubscript{$\infty$}

régime transitoire

Pour la \textbf{réponse à un échelon}, la précision est définie par
\textbf{l'écart statique} \textbf{ou écart de position}
\(\mathbf{\epsilon}_{\mathbf{s}}\) \textbf{(ou}
\(\mathbf{\epsilon}_{\mathbf{p}}\)\textbf{)} . C'est le cas des systèmes
régulateurs.

\textbf{Réponse à une rampe}

\textbf{$\varepsilon$\textsubscript{r}}

\textbf{Ecart de traînage ou écart de suivi $\varepsilon$\textsubscript{r}}

s(t)

0

t

régime transitoire

Pour la \textbf{réponse à une rampe}, la précision est définie par
\textbf{l'écart de traînage ou l'écart de suivi}
\(\mathbf{\epsilon}_{\mathbf{r}}\) \textbf{(ou}
\(\mathbf{\epsilon}_{\mathbf{t}}\)\textbf{)}. C'est le cas des systèmes
suiveurs.

À condition que le système asservi soit \textbf{stable}, la
\textbf{précision} n'est définie que pour un système ayant des
\textbf{grandeurs d'entrée et de sortie} \textbf{de même nature}.

L'écart : \(\epsilon(t) = e(t) - s(t)\)

L'écart statique~:
\(\epsilon_{s} = \lim_{t \rightarrow \infty}{\epsilon(t)}\)

Si le système est stable et que l'entrée et la sortie sont de même
nature, la \textbf{précision} est caractérisée par~:

\begin{itemize}
\item
  \textbf{l\textquotesingle écart statique~absolu (dans l'unité de la
  consigne)}: \(\epsilon_{s} = E_{0} - s_{\infty}\)
\item
  ou par \textbf{l\textquotesingle écart statique relatif (en \%)~}:
\end{itemize}

\(\epsilon_{s\%}\  = \left| \frac{\epsilon_{s}}{\Delta E} \right|\) pour
une \ul{consigne en échelon d'amplitude} \(E_{0}\)\ul{.}

Si l'écart statique est \textbf{nul}, on dira que le système est
\textbf{précis}.

\[t_{init}\]

\[t\]

\[s_{\infty}\]

\[\Delta S\]

\[\Delta E\]

\[s(t)\ e(t)\]

\[\epsilon_{s}\]

\[E_{0}\]

\begin{longtable}[]{@{}
  >{\raggedright\arraybackslash}p{(\columnwidth - 2\tabcolsep) * \real{0.1227}}
  >{\raggedright\arraybackslash}p{(\columnwidth - 2\tabcolsep) * \real{0.8773}}@{}}
\toprule\noalign{}
\begin{minipage}[b]{\linewidth}\raggedright
\begin{quote}
\SAimage{images/image18.png}{13mm}{14mm}
\end{quote}
\end{minipage} & \begin{minipage}[b]{\linewidth}\raggedright
\textbf{Robot chirurgical Da vinci}

\SAimage{images/image29.png}{68mm}{38mm}

x(t) {[}mm{]}

\SAimage{images/image30.pdf}{90mm}{75mm}

t {[}ms{]}

\textbf{Déterminer l\textquotesingle écart statique et
l\textquotesingle écart statique relatif.}

\begin{quote}
Les grandeurs d'entrée et de sortie sont de même nature.

\[\epsilon_{s} = E_{0} - x_{\infty} = 7 - 8 = - 1\ mm\]

\[\epsilon_{s\%} = \left| \frac{\epsilon_{s}}{\Delta E} \right| = \left| \frac{- 1}{7 - 5} \right| = 0,5 = \frac{50}{100} = 50\%\]
\end{quote}

\textbf{Déterminer l\textquotesingle écart de trainage.}

x(t) {[}mm{]}

t {[}ms{]}

\SAimage{images/image38.png}{90mm}{66mm}
\end{minipage} \\
\midrule\noalign{}
\endhead
\bottomrule\noalign{}
\endlastfoot
\end{longtable}

\subsection{Dépassement}\label{duxe9passement}

Une réponse présente un \textbf{dépassement} si, pour une entrée en
échelon, la grandeur de sortie dépasse sa \textbf{valeur finale}. La
courbe de sortie présente alors un extrémum. Ce dépassement \textbf{est
évalué à partir de la réponse et de sa valeur finale.} Pour certains
systèmes, il est impératif qu'il n'y ait aucun dépassement.

\[\Delta D\]

\[s_{0}\]

\[t\]

\[t_{1}\]

\[t_{2}\]

\[s_{\infty}\]

\[D_{2}\]

\[s_{\max}\]

\[\Delta S\]

\[s(t)\]

\[D_{1}\]

Moins il y a d\textquotesingle oscillations, plus le système est
\textbf{amorti}. Dans certaines applications elles seront à éviter
(positionnement d'une machine à commande numérique, ...). La valeur du
premier dépassement (le plus critique) se détermine à partir de la
réponse indicielle.

\[D_{1}(\%) = 100.\left| \frac{\Delta D}{\Delta S} \right| = 100.\left| \frac{s_{\max} - s_{\infty}}{s_{\infty} - s_{0}} \right|\]

\textbf{L'amplitude du premier} \textbf{dépassement}, notée
\emph{\textbf{D\textsubscript{1}}}, par rapport à la valeur finale:

\[D_{1}(\%) = 100.\left| \frac{\Delta D}{\Delta S} \right| = 100.\left| \frac{s_{\max} - s_{\infty}}{s_{\infty} - s_{0}} \right|\]

\subsubsection{}\label{section-1}

\subsubsection{ Le dépassement n'existe que pour des systèmes d'ordre 2 ou supérieur. Un système d'ordre 1 ne peut pas avoir de dépassement.}\label{le-duxe9passement-nexiste-que-pour-des-systuxe8mes-dordre-2-ou-supuxe9rieur.-un-systuxe8me-dordre-1-ne-peut-pas-avoir-de-duxe9passement.}

Ce ne sont pas les dépassements par rapport à la valeur en entrée mais
bien par rapport à la valeur finale de la sortie.

\begin{longtable}[]{@{}
  >{\raggedright\arraybackslash}p{(\columnwidth - 2\tabcolsep) * \real{0.1227}}
  >{\raggedright\arraybackslash}p{(\columnwidth - 2\tabcolsep) * \real{0.8773}}@{}}
\toprule\noalign{}
\begin{minipage}[b]{\linewidth}\raggedright
\begin{quote}
\SAimage{images/image18.png}{13mm}{14mm}
\end{quote}
\end{minipage} & \begin{minipage}[b]{\linewidth}\raggedright
\textbf{Robot chirurgical Da vinci}

\SAimage{images/image29.png}{68mm}{38mm}Le
robot Da Vinci est un robot médical dirigé par un chirurgien pour
réaliser des opérations, principalement au niveau de l'abdomen.
L'opération se déroule à distance à l'aide d'un écran et de dispositifs
haptiques. Les mouvements sont tous asservis et requièrent un réglage
fin.

\textbf{Déterminer la valeur du premier dépassement.}

x(t) {[}mm{]}

\SAimage{images/image30.pdf}{90mm}{75mm}

t {[}ms{]}

Le système est stable car pour une entrée en échelon, la sortie converge
vers une valeur constante (EBSB).

\[D_{1} = \left| x_{\max} - x_{\infty} \right| = |8,9 - 8| = 0,9\ mm\]

D\textsubscript{1\%}=\(\left| \frac{x_{\max} - x_{\infty}}{x_{\infty} - x_{0}} \right| = \left| \frac{0,9}{8 - 5} \right| = \ 0,3 = 30\%\)
\end{minipage} \\
\midrule\noalign{}
\endhead
\bottomrule\noalign{}
\endlastfoot
\end{longtable}

\begin{longtable}[]{@{}
  >{\raggedright\arraybackslash}p{(\columnwidth - 2\tabcolsep) * \real{0.1227}}
  >{\raggedright\arraybackslash}p{(\columnwidth - 2\tabcolsep) * \real{0.8773}}@{}}
\toprule\noalign{}
\begin{minipage}[b]{\linewidth}\raggedright
\begin{quote}
\SAimage{images/image18.png}{13mm}{14mm}
\end{quote}
\end{minipage} & \begin{minipage}[b]{\linewidth}\raggedright
\textbf{Stabilisateur}

Un essai indiciel sur le stabilisateur pour un déplacement de 2$^\circ$ est
donné ci-dessous.

\SAimage{images/image33.png}{90mm}{74mm}

\textbf{Déterminer la valeur du premier dépassement.}
\end{minipage} \\
\midrule\noalign{}
\endhead
\bottomrule\noalign{}
\endlastfoot
\end{longtable}

\section{Sources}\label{sources}

Ce cours a été élaboré à l'aide de nombreuses ressources~provenant de
différents collègues~de l'UPSTI.
